\section{AI in Meeting Support Systems (RQ2)}

In addition to the fundamentals of Group Support Systems (GSS) outlined in
\autoref{sec:gss-collaboration-tools}, modern approaches
specifically designed for meeting support systems can be identified, some of which are supplemented by
functions that are more or less supported by AI. This chapter analyzes these systems to identify potential
gaps and to serve as a source of inspiration for the development of the prototype.

% ---------------------------------------------------------------------------------------------------------------------
\subsection{Overview of existing Meeting AI Tools}
\label{sec:existing-meeting-ai-tools}
There is a wide variety of AI-based tools that provide support for meetings. Most of these tools focus on
transcribing meeting content. Well-known examples include fireflies.ai \cite{Firefliesai1AI}, Otter.ai
\cite{otteraiOtterMeetingAgent}, and Jamie \cite{jamieaiJamieYourPersonal}.
AI transcription features are also integrated into established meeting platforms such as Google Meet
\cite{workspaceGoogleMeetTelefon}, Zoom \cite{zoomKINotizentoolIhrKIMeetingAssistent},
and Microsoft Teams. \cite{microsoftVideokonferenzenBesprechungenAnrufe}
More advanced features include the automatic extraction of action items. Depending on the tool, these
features are available in real time or only after the meeting has ended.
Microsoft Teams specifically offers a feature called “Facilitator,” \cite{microsoftSetFacilitatorMicrosoft}
which provides additional capabilities beyond the features already mentioned (transcription, identification
of action items): For example, it can edit the agenda if instructed to do so by a user with the appropriate
permissions. It also displays a timeline to track topics and the flow of the meeting, and sends reminders to
manage time. (Once halfway through the meeting and again shortly before it ends.) These interactions take
place via the integrated chat.
Google Meet, for example, also offers a live speech translation feature. \cite{googleLearnSpeechTranslation}

Nevertheless, contemporary meeting platforms remain primarily technical infrastructures for conducting
meetings. Their design does not systematically account for the professional knowledge, needs, and
methodological repertoire of facilitators \cite{hughesKeeperSynchronousOnline2021}, nor are they conceived to
support the range of tools, methods, practices, and techniques on which facilitation typically relies.

There are also platforms specialized for facilitation needs, such as Stormz \cite{GetStartedStormz}, offer AI
support across different phases of a meeting, including brainstorming.
\cite{aitinkerersclubStormzAIFeature2023}

Platforms such as SessionLabs curate methods from the facilitation community and, with the help of an AI
assistant, make it possible to find suitable methods and generate an agenda tailored to the meeting's
objectives.\cite{sessionlabsDesignFasterAI2025}

% ---------------------------------------------------------------------------------------------------------------------
\subsection{Related Research Approaches to AI-Supported Meeting-Systems}
\label{sec:research-approaches-to-ai-supported-meeting-systems}
Several works in the literature pursue a comparable approach or related research questions . Among them, two
papers proved particularly relevant to this thesis. Their key findings are summarized below to clarify how
they inform the subsequent work.

\textbf{Are We On Track?}

In an empirical study involving 15 knowledge workers, Chen et al. examined how AI can
support goal orientation and intentionality in meetings through passive visual cues and active, inquiry-based
reflection prompts. \cite[p. 1]{chenAreWeTrack2025}
To increase ecological validity, the authors used past meeting recordings provided by the study participants
and applied a video-stimulated method in which the participants
rewatched the same meeting situation, with the addition that it now included AI interventions.
They evaluated counter-scenarios with simulated AI interventions. \cite[p. 4]{chenAreWeTrack2025}
Technically, the prototypes developed were based primarily on GPT-4, which processed transcripts in simulated
real time to passively visualize topic progression using D3.js or actively display context-related reflection questions.
\cite[p. 6]{chenAreWeTrack2025}
The results show that passive visualizations maintain the flow of conversation and promote individual
understanding, whereas active AI questions, while triggering immediate adjustments, carry the risk of
disruptions to the discourse. \cite[p. 2]{chenAreWeTrack2025} However, due to the passive designs,
particularly in the visualization system,it was found that participants overlooked or ignored them. \cite[p.
2]{chenAreWeTrack2025}
Based on the participants’ responses, the authors conclude that it should be possible to
dynamically adjust the intensity of the intervention to the subjective need for support in order to ensure a
balance between
democratic participation, efficiency, and user control. \cite[p. 13]{chenAreWeTrack2025}

One limitation of this study is that the system was not tested in the context of a live meeting
.

\textbf{Observe, Ask Intervene: designing AI Agents for More Inclusive Meetings}

In the study by Houtti et al., the authors examined the design of an AI agent as a so-called “virtual
co-host” to specifically promote inclusion and equal participation in video conferences.\cite
{houttiObserveAskIntervene2025}
First, ideation sessions were conducted in which participants were asked to design a virtual co-host
that could intervene to balance out inequalities in participation. \cite[p.
6]{houttiObserveAskIntervene2025}

This concept was then evaluated in an exploratory laboratory study involving 68 participants in 12
groups, using a structured decision-making task (“Lost at Sea”) on Discord, followed by questionnaires
on meeting quality and 18 qualitative interviews. \cite[p. 14]{houttiObserveAskIntervene2025}

Technically, the prototype was implemented as a Discord bot in Python (`Discord.py`), which used Selenium and
custom JavaScript to track the cumulative speaking times of all participants in the browser interface in
real time. The bot communicated exclusively via private direct messages, generated
individual bar charts showing speaking time, and always waited for a speaking pause of
at least five seconds before sending messages to avoid interrupting the conversation.
\cite[p.6]{houttiObserveAskIntervene2025}

A tendency was observed to design the AI facilitator in a rather reactive manner. The participants wanted to retain a
certain degree of control over the co-host. \cite[p.6]{houttiObserveAskIntervene2025} Nevertheless, the
authors concluded that this is not always feasible, because they suspected that a system that
must always be triggered might remain unused.
Based on this, the authors derive a framework they call “Observe, Ask, Intervene.”
The AI agent continuously observes the meeting but asks a human participant for permission before intervening.
Nearly all interviewees stated that they would prefer to be asked before an intervention affecting them takes
place. \cite[p.15] {houttiObserveAskIntervene2025}

Specifically, this worked as follows: After certain thresholds for distributed speaking time were reached,
the bot sent a private direct message with questions to people who had a low level of participation. Only
after the person responded accordingly did the agent send the speaking time charts.

Although private feedback was perceived as socially acceptable\cite[p.17]{houttiObserveAskIntervene2025},
the bot did not exert sufficient social or normative pressure, which is why facilitators or dominant speakers
largely ignored or rationalized the AI’s critical feedback.
\cite[p.18]{houttiObserveAskIntervene2025}

It should also be noted that the system is a rule-based prototype that does not utilize a generative LLM pipeline.

% ---------------------------------------------------------------------------------------------------------------------
\subsection{Proactive AI-Interactions: From Turn-Taking to intrinsic Motivation with the Inner Thoughts Framework}
\label{sec:proactive-ai-interactions}

So an AI-System can proactivly interact and therefor reduese the need for explicit prompts by the meeting participants
it should be able to proactively recognize when it is its turn in an interaction sequence.

In dyadic human-AIinteractions, pauses or so-called “stop words” are often
used for this purpose to identify a change in speaker. In multi-party scenarios (e.g., group conversations),
however, this becomes significantly more complex and leads to the “multi-party turn-taking
problem”. \cite{sacksSimplestSystematicsOrganization1974}

Some approaches model the problem as a machine learning task by attempting to predict the next speaker based
on previous speaker changes.
\cite{bayserLearningMultiPartyTurnTaking2019}
Bayer et al. employ traditional algorithms such as Maximum Likelihood Estimation (MLE) and Support Vector
Machines (SVMs), as well as deep learning architectures including Convolutional Neural Networks (CNN) and
Long Short-Term Memory (LSTM) networks.

With these approaches, however, it remains unclear why an agent should interact at a specific point in time.
Due to the design of the ML algorithms, the timing of the interaction is decoupled from its cause or is only
implicitly included in the model, but cannot be explicitly explained.

A naive approach to turn-taking prediction would be to prompt an LLM to ask who is next.

Liu et al. investigated various language models (including GPT-3.5, GPT-4o, and fine-tuned variants) with
regard to their ability to predict the next speaker (next-speaker prediction). They showed that these
approaches work well when speaker identification (turn allocation) is explicit.
\cite[p.4]{liuProactiveConversationalAgents2025}

This makes sense, since the current speaker explicitly assigns the next turn to a specific person, and the
conversational context thus contains clear semantic cues, such as direct mentions of names—that an LLM should
be able to process well.

However, in self-selection scenarios, where a participant chose themselves as the next speaker without being
asked to do so, all models performed no better than chance.
\cite[p.4]{liuProactiveConversationalAgents2025}

The authors conclude from this that the conversational context alone is insufficient for self-selection*
scenarios, and that hidden, internal factors (such as interest, relevance, or motivation to engage) play a
decisive role. \cite[p.4]{liuProactiveConversationalAgents2025} Based on a
thinking-aloud study with 24 participants, they concluded that people are driven by intrinsic motivation to
voice their thoughts. \cite[p.2]{liuProactiveConversationalAgents2025}

Building on this, they modeled a framework for chatbots that they call “Inner Thoughts.” An LLM-based agent
continuously generates internal thoughts in parallel with the conversation and evaluates them based on a
defined heuristic (a fixed prompt in the LLM call). If the intrinsic motivation (i.e., the evaluation of a
thought) exceeds a configurable threshold, the agent decides whether to participate, taking into account the
turn-taking type. \cite{liuProactiveConversationalAgents2025}

In chat evaluations, this system achieved significantly higher performance in self-selection scenarios
than conventional baselines, and human evaluators perceived the conversations as significantly
more natural, intelligent, and coherent. \cite[p.2]{liuProactiveConversationalAgents2025}
